% Bản nháp chương 3 theo dàn bài mới, viết với bộ mẫu ở style/mau/.
% ===========================================================================
\chapter{Tối ưu tới mức nào là đủ}
\label{sec:enough}

Mọi biểu đồ hội tụ trong báo cáo đo bằng $f(w_k) - \fstar$, còn cửa hàng định
giá thì chỉ quan tâm sai số giá. Hai đại lượng ấy chưa gặp nhau ở chỗ nào, nên
câu hỏi hiển nhiên nhất vẫn còn treo: giải bài toán tối ưu hóa tới mức nào thì
lời giải ngừng tốt lên đối với người dùng nó. Chương này trả lời bằng một thí
nghiệm riêng, và câu trả lời ấy đặt ra ngưỡng dừng cho toàn bộ phần còn lại.

\section{Cách đo}
\label{sec:enough-method}

Nhóm ghi RMSE trên tập kiểm tra tại từng điểm ghi lịch sử của từng lần chạy, cho
cấu hình tốt nhất của bảy nhóm thí nghiệm. Phép đo này nằm trong cửa sổ mà đồng
hồ đã dừng, cùng chỗ với phép tính $f(w_k)$ và chuẩn gradient, nên chi phí của
nó không bị tính vào thời gian chạy; bộ đếm số lần đánh giá hàm cũng được khôi
phục sau mỗi lần ghi. Chi tiết cài đặt nằm ở mục~\ref{sec:recorder}.

Loại trừ chi phí ấy vẫn chưa đủ để thời gian đo được giữ nguyên. Mỗi lần ghi,
phép tính RMSE quét toàn bộ ma trận kiểm tra \SI{85}{\mebi\byte} xen vào giữa
các vòng lặp và đẩy ma trận huấn luyện \SI{342}{\mebi\byte} ra khỏi cache, nên
các vòng lặp ngay sau đó chạy chậm hơn. Đo trực tiếp trên gradient descent, bật
monitor làm thời gian chạy tăng 13\%, từ \SI{5.72}{\second} lên
\SI{6.47}{\second}. Vì thế chương này chỉ dùng số vòng lặp, vốn không bị ảnh
hưởng, còn mọi con số thời gian đều lấy từ chương~\ref{sec:compare}, nơi thí
nghiệm chạy không có monitor.

Cụ thể, sáu cấu hình hội tụ cho 153 cặp giá trị gồm sai số tối ưu hóa và RMSE
tương ứng. SGD với độ dài bước hằng không hội tụ nên nó nằm ngoài phần dựng
ngưỡng và được xét riêng ở mục~\ref{sec:enough-sgd}.

Mốc để so là RMSE tại nghiệm đóng $\wstar$, bằng \num{0.2060141}, ứng với sai số
giá \num{4245} đơn vị tiền trên một chiếc máy giá trung vị theo
\eqref{eq:rmse-to-price}. Quanh mức ấy, một đơn vị RMSE đổi được khoảng
\num{22800} đơn vị tiền, nên chênh lệch RMSE ở chữ số thập phân thứ năm đã nằm
dưới một đồng và không còn ý nghĩa với người bán máy.

\section{RMSE ngừng cải thiện từ đâu}
\label{sec:enough-curve}

\resultfig{rmse_vs_gap}
  {RMSE trên tập kiểm tra theo sai số tối ưu hóa tại chính điểm đo. Trục hoành
   chạy từ sai số lớn bên trái sang sai số nhỏ bên phải, tức theo chiều các
   thuật toán tiến. Trục tung cắt sát mức của nghiệm đóng nên các đường đi vào
   khung từ mép trên.}
  {fig:rmse-vs-gap}

Hình~\ref{fig:rmse-vs-gap} cho thấy cả sáu đường đều phẳng lại và nằm chồng lên
mức của nghiệm đóng, nhưng chúng phẳng ở những chỗ khác nhau. Ở cùng một mức sai
số tối ưu hóa, các thuật toán không cho cùng một RMSE, vì $f(w_k) - \fstar$ chỉ
là một số vô hướng còn hai điểm có cùng giá trị hàm mục tiêu vẫn nằm ở hai chỗ
khác nhau trong không gian tham số. Ngưỡng cần tìm vì thế phải lấy theo trường
hợp xấu nhất trong cả sáu đường, không lấy theo một đường riêng.

Bảng~\ref{tab:app-threshold} ghi kết quả gộp: với mỗi mức sai số tối ưu hóa, cột
thứ hai là sai lệch giá lớn nhất trong mọi điểm ghi có sai số dưới mức đó.

\begin{table}[tp]
  \centering
  \caption{Sai lệch giá lớn nhất so với nghiệm đóng, gộp 153 điểm ghi của sáu
    cấu hình hội tụ, tính trên một chiếc máy giá trung vị \num{18555}.}
  \label{tab:app-threshold}
  \begin{tabular}{lr}
    \toprule
    Sai số tối ưu hóa dưới & Sai lệch giá lớn nhất \\
    \midrule
    $10^{-2}$ & \num{634.4} \\
    $10^{-3}$ & \num{108.8} \\
    $10^{-4}$ & \num{12.4}  \\
    $10^{-5}$ & \num{1.56}  \\
    $10^{-6}$ & \num{0.61}  \\
    $10^{-7}$ & \num{0.12}  \\
    $10^{-8}$ & \num{0.02}  \\
    \bottomrule
  \end{tabular}
\end{table}

Sai lệch giá giảm chậm hơn sai số tối ưu hóa. Khớp một hàm lũy thừa vào 85 điểm
ghi trải từ $10^{-12}$ tới $10^{-1}$ cho số mũ \num{0.68}, tức cắt sai số tối ưu
hóa đi 10 lần chỉ giảm được sai lệch giá khoảng 5 lần. Hai dự đoán đơn giản đều
không đúng hẳn: vì $f$ là hàm bậc hai nên $f(w_k) - \fstar$ tỉ lệ với
$\norm{w_k - \wstar}^{2}$, do đó số mũ phải bằng \num{0.5} nếu sai lệch giá tỉ lệ
với $\norm{w_k - \wstar}$ và bằng 1 nếu nó tỉ lệ với bình phương khoảng cách ấy.
Giá trị đo được nằm giữa hai mức, nên không quy về một trong hai bức tranh đó
được. Điều dùng được là hệ quả thực dụng: mỗi chữ số chính xác thêm về phía tối
ưu hóa mua được ngày càng ít về phía bài toán ứng dụng.

\section{Ngưỡng đủ dùng}
\label{sec:enough-threshold}

Đặt $\eps_{\text{app}}$ là sai số tối ưu hóa lớn nhất mà từ đó trở xuống, mọi
điểm ghi đều cho sai lệch giá nằm trong một dung sai định trước. Bảng~\ref{tab:app-threshold} cho ngay
hai giá trị. Với dung sai một đơn vị tiền trên máy giá trung vị,
$\eps_{\text{app}} = 1{,}35 \cdot 10^{-6}$. Với dung sai mười đơn vị, tức khoảng
hai phần nghìn của khoản sai số mà bản thân mô hình vốn mắc phải,
$\eps_{\text{app}} = 4{,}0 \cdot 10^{-5}$.

Con số thứ nhất đáng chú ý vì một lý do ngoài dự kiến. Ngưỡng $10^{-6}$ mà báo
cáo dùng từ đầu là một quy ước, chọn vì nó tròn và vì nó nằm đủ xa giới hạn số
học; hóa ra nó chỉ chặt hơn ngưỡng có căn cứ từ bài toán ứng dụng 1,35 lần. Quy
ước ấy vì thế không cần đổi, nhưng từ chương này nó có một lý do thay cho thói
quen.

\begin{table}[tp]
  \centering
  \caption{Số vòng lặp cần để đạt hai ngưỡng. Cột giữa ứng với dung sai mười đơn
    vị tiền, cột phải ứng với một đơn vị. Với SGD, một vòng lặp ở đây là một lượt
    duyệt dữ liệu.}
  \label{tab:two-thresholds}
  \begin{tabular}{lrr}
    \toprule
    Thuật toán
      & $\eps_{\text{app}} = 4{,}0 \cdot 10^{-5}$
      & $\eps_{\text{app}} = 1{,}35 \cdot 10^{-6}$ \\
    \midrule
    Newton ($t = 1$)         & 1  & 1  \\
    L-BFGS ($m = 20$)        & 6  & 8  \\
    GD ($t = 2/L$)           & 10 & 20 \\
    GD backtracking          & 10 & 15 \\
    AGD + restart            & 15 & 20 \\
    SGD bậc thang ($B = 64$) & 21 & không đạt \\
    \bottomrule
  \end{tabular}
\end{table}

Nới dung sai từ một đơn vị lên mười đơn vị cắt số vòng lặp của các phương pháp
bậc một xuống khoảng một nửa, chẳng hạn gradient descent đi từ 20 vòng còn 10
(bảng~\ref{tab:two-thresholds}). Nhưng thay đổi đáng kể hơn nằm ở dòng cuối: SGD
bậc thang đạt ngưỡng nới sau 21 lượt duyệt dữ liệu, trong khi ở ngưỡng chặt nó
không đạt trong toàn bộ ngân sách 30 lượt. Nói cách khác, chính dung sai đặt ra ở
chương này quyết định SGD có được coi là phương pháp dùng được hay không, chứ
không phải bản thân thuật toán.

\section{Hai quan sát ngoài lề}
\label{sec:enough-sgd}

Chạy tiếp sau $\wstar$ không cải thiện thêm, và có lúc còn hơi quá đà. Gradient
descent với $t = 2/L$ dừng ở sai số $4{,}9 \cdot 10^{-8}$ và cho RMSE
\num{0.2060116}, tức thấp hơn mức của nghiệm đóng \num{0.06} đơn vị tiền. Điều
này không mâu thuẫn với \eqref{eq:closed-form}: $\wstar$ cực tiểu hàm mục tiêu
trên tập huấn luyện chứ không cực tiểu sai số trên tập kiểm tra, nên một điểm
dừng sớm vẫn có thể nằm hơi tốt hơn về mặt dự báo. Mức chênh ở đây quá nhỏ để
khai thác, nhưng nó là dạng thu gọn của hiện tượng dừng sớm quen thuộc trong học
máy.

Chỉ duy nhất SGD với độ dài bước hằng là ngoại lệ thật sự. Nó dừng ở sai số
$8{,}2 \cdot 10^{-4}$, không phải vì hết ngân sách mà vì bước hằng chỉ đưa nghiệm
tới một lân cận quanh $\wstar$ rồi dao động trong đó, cơ chế được trình bày ở
chương~\ref{sec:mechanisms}. Sai số ấy nằm trên cả hai ngưỡng, và hệ quả đo được
là RMSE \num{0.20986} so với \num{0.20601}, tức đắt thêm \num{87.7} đơn vị tiền
trên mỗi chiếc máy giá trung vị. Đây là cấu hình duy nhất trong báo cáo mà lựa
chọn thuật toán đổi được con số cửa hàng đưa ra cho khách.

\section{Hệ quả cho phần còn lại}
\label{sec:enough-consequence}

Từ chương này trở đi, một thuật toán được gọi là đủ tốt khi nó đưa sai số xuống
dưới $10^{-6}$, và con số ấy giờ có nghĩa cụ thể: sai lệch giá dưới một đơn vị
tiền trên một chiếc máy giá trung vị, tức dưới hai phần vạn của khoản sai số
\num{4245} mà bản thân mô hình vốn đã mắc phải.

Điều đó cũng đặt lại tỉ lệ cho mọi so sánh về sau. Khoảng cách giữa mô hình tốt
nhất và mô hình tệ nhất trong báo cáo, xét theo chất lượng định giá, là 87,7 đơn
vị tiền của SGD bước hằng, trong khi bản thân mô hình sai trung bình \num{4245}
đơn vị. Nói cách khác, toàn bộ phần tối ưu hóa quyết định khoảng 2\% sai số cuối
cùng, còn 98\% còn lại thuộc về việc chọn mô hình và chuẩn bị dữ liệu.
Chương~\ref{sec:compare} so sánh các thuật toán trong đúng khung tỉ lệ ấy: không
phải để tìm thuật toán dự đoán giá đúng hơn, mà để tìm thuật toán tới đích rẻ
hơn.
